\section{Taller resuelto: Dimensionar combustible por cierre de misión}

\noindent\textbf{Nivel:} Intermedio.

\apartado{1. Enunciado}{Una fase de reactor debe recorrer 1400 km en calma a 7000 m ISA y CL de máximo alcance. La masa final, que incluye carga y reserva, es 6500 kg. Con S=42 m², polar 0.024+0.044CL² y cT=1.7e-5 kg/(N s), determine masa inicial y combustible de la fase.}

\apartado{2. Datos}{El combustible de la fase es desconocido y forma parte de masa inicial.}

\apartado{3. Incógnitas}{Las magnitudes y comprobaciones solicitadas en el caso.}

\apartado{4. Sistema de referencia}{Crucero horizontal y TAS variable para CL fijo.}

\apartado{5. Hipótesis físicas}{Polar, densidad y consumo constantes; no hay ascenso ni transición dentro de la fase.}

\apartado{6. Modelo matemático}{Dimensionar combustible por cierre de misión}

\apartado{7. Ecuaciones simbólicas}{Se parte de los balances y definiciones siguientes.}

\begin{equation}R=K_m(\sqrt{m_i}-\sqrt{m_f})\end{equation}

\begin{equation}m_i=(\sqrt{m_f}+R/K_m)^2\end{equation}

\begin{equation}m_{comb}=m_i-m_f\end{equation}

\apartado{8. Desarrollo algebraico}{Las operaciones se realizan en este orden, conservando las condiciones del modelo.}

\begin{enumerate}

\item Se calcula CL de máximo alcance y el coeficiente Km de la ley.

\item Se despeja masa inicial antes de restar combustible; no se utiliza una masa media arbitraria.

\item Se verifican velocidades extremas y que el alcance se reconstruya con la masa hallada.

\end{enumerate}

\apartado{9. Sustitución numérica}{Se conservan los resultados intermedios sin redondear y se muestran todas las operaciones escalares disponibles.}

\textbf{Desarrollo de las magnitudes intermedias.} De $dm/dt=-c_TD=-c_Tmg_0/(L/D)$, separar variables da $E=(L/D)\ln(m_i/m_f)/(g_0c_T)$. Con CL fijo, $V=\sqrt{2g_0m/(\rho SC_L)}$; sustituir en $dx=Vdt$ deja integrar $m^{-1/2}$, cuya primitiva es $2\sqrt m$. Por eso $R=K_m(\sqrt{m_i}-\sqrt{m_f})$, con $K_m=2(L/D)\sqrt{2g_0/(\rho SC_L)}/(g_0c_T)$. Para reactor a CL constante, máximo alcance se obtiene maximizando $\sqrt{C_L}/C_D$, que da $C_L=\sqrt{C_{D0}/(3k)}$. Despejar la masa inicial requiere primero $\sqrt{m_i}=\sqrt{m_f}+R/K_m$ y después elevar al cuadrado.\begin{enumerate}

\item Elegir el óptimo de alcance de reactor.

\[\fitmath{C_L=\sqrt{\frac{0.024}{3\cdot 0.044}}\approx 0.4264014\,}\]

\item Evaluar resistencia adimensional.

\[\fitmath{C_D=0.024+0.044\cdot {C_L}^{2}}\]

\[\fitmath{C_D=0.024+0.044\cdot {(0.4264014327)}^{2}\approx 0.032\,}\]

\item Obtener eficiencia de esa condición.

\[\fitmath{\mathrm{LD}=\frac{C_L}{C_D}}\]

\[\fitmath{\mathrm{LD}=\frac{(0.4264014327)}{(0.032)}\approx 13.32504\,}\]

\item Calcular el coeficiente completo de la integral.

\[\fitmath{\mathrm{Km}=\frac{2\cdot \mathrm{LD}}{9.80665\cdot 1.7\times10^{-5}}\cdot \sqrt{\frac{2\cdot 9.80665}{0.589500721\cdot 42\cdot C_L}}}\]

\[\fitmath{\mathrm{Km}=\frac{2\cdot (13.32504477)}{9.80665\cdot 1.7\times10^{-5}}\cdot \sqrt{\frac{2\cdot 9.80665}{0.589500721\cdot 42\cdot (0.4264014327)}}\approx 217885.5\,\mathrm{m/kg^{.5}}}\]

\item Cerrar el alcance prescrito antes de estimar combustible.

\[\fitmath{\mathrm{mi}={\left(\sqrt{6500}+\frac{1400000}{\mathrm{Km}}\right)}^{2}}\]

\[\fitmath{\mathrm{mi}={\left(\sqrt{6500}+\frac{1400000}{(217885.4692)}\right)}^{2}\approx 7577.349\,\mathrm{kg}}\]

\item Restar masa final que ya incluye reserva.

\[\fitmath{\mathrm{fuel}=\mathrm{mi}-6500}\]

\[\fitmath{\mathrm{fuel}=(7577.34935)-6500\approx 1077.349\,\mathrm{kg}}\]

\item Comprobar velocidad inicial.

\[\fitmath{\mathrm{Vi}=\sqrt{\frac{2\cdot \mathrm{mi}\cdot 9.80665}{0.589500721\cdot 42\cdot C_L}}}\]

\[\fitmath{\mathrm{Vi}=\sqrt{\frac{2\cdot (7577.34935)\cdot 9.80665}{0.589500721\cdot 42\cdot (0.4264014327)}}\approx 118.6473\,\mathrm{m/s}}\]

\item Comprobar velocidad final.

\[\fitmath{\mathrm{Vf}=\sqrt{\frac{2\cdot 6500\cdot 9.80665}{0.589500721\cdot 42\cdot C_L}}}\]

\[\fitmath{\mathrm{Vf}=\sqrt{\frac{2\cdot 6500\cdot 9.80665}{0.589500721\cdot 42\cdot (0.4264014327)}}\approx 109.8894\,\mathrm{m/s}}\]

\end{enumerate}





\apartado{10. Resultados con unidades}{Los resultados pertenecen al modelo y al intervalo de datos declarados.}

\begin{center}\begin{tabular}{lr}\toprule Magnitud&Resultado\\\midrule

Masa inicial & $7577.349\,\mathrm{kg}$\\

Combustible de fase & $1077.349\,\mathrm{kg}$\\

TAS inicial & $118.6473\,\mathrm{m\,s^{-1}}$\\

TAS final & $109.8894\,\mathrm{m\,s^{-1}}$\\

\bottomrule\end{tabular}\end{center}

\apartado{11. Interpretación}{La masa final ya contiene la reserva; el combustible calculado es sólo el quemado en esta fase.}

\apartado{12. Comprobación independiente y alcance}{Resolver por raíz numérica la diferencia entre alcance calculado y prescrito; comprobar cierre de masa.}

\apartado{13. Limitaciones}{La factibilidad necesita capacidad de tanque, masa máxima, empuje y límites de Mach, que no se han dado.}

\par\smallskip\noindent{\footnotesize\textit{Datos y solución originales. Las sustituciones registradas se reproducen con SymPy y Studium; eso no establece validación experimental ni revisión experta.}}
